\documentclass[11pt,a4paper]{article}
\usepackage[a4paper,top=15mm,bottom=15mm,left=16mm,right=16mm]{geometry}
\usepackage[T1]{fontenc}
\usepackage{amsmath,amssymb}
\usepackage{enumitem}
\usepackage{array}
\usepackage{xcolor}

\setlength{\parindent}{0pt}
\setlength{\parskip}{4pt}
\renewcommand{\arraystretch}{1.2}
\setlist[itemize]{leftmargin=*,itemsep=2pt,topsep=2pt}

\begin{document}

\begin{center}
  {\Large\bfseries Worked Solutions}\\[2pt]
  {\large Picture Sequences and the \boldmath\(n\)th Term}
\end{center}

\section*{Part 1: Linear sequences}

\paragraph{Worked example: row of squares.}
Each pattern has 1 matchstick at the start, and then \(n\) groups of 3 matchsticks. Hence pattern \(n\) has
\[
1+3n=3n+1
\]
matchsticks.

\[
\begin{array}{@{}lccccccc@{}}
\text{Pattern }n&1&2&3&4&5&10&n\\ \hline
\text{Matchsticks}&4&7&10&13&16&31&3n+1
\end{array}
\]

\paragraph{Q1. Row of houses.}
Each pattern has 1 starting matchstick, and then \(n\) groups of 5 matchsticks. Therefore the number of matchsticks is
\[
1+5n=5n+1.
\]
The completed table is:

\[
\begin{array}{@{}lccccccc@{}}
\text{Pattern }n&1&2&3&4&5&10&n\\ \hline
\text{Matchsticks}&6&11&16&21&26&51&5n+1
\end{array}
\]

\paragraph{Q2. Tables in a caf\'e.}
Each table has 2 chairs, one above and one below, giving \(2n\) chairs. There is also 1 chair at each end, so the total is
\[
2n+2.
\]
The completed table is:

\[
\begin{array}{@{}lccccccc@{}}
\text{Pattern }n&1&2&3&4&5&10&n\\ \hline
\text{Chairs}&4&6&8&10&12&22&2n+2
\end{array}
\]
The rings should show \(n\) pairs of chairs plus the two end chairs.

\paragraph{Q3. Dot sequences.}
\begin{enumerate}[label=(\alph*),leftmargin=*,itemsep=4pt]
\item A cross has one centre dot and four arms of \(n\) dots, so
\[
1+4n=4n+1.
\]
\[
\begin{array}{@{}lccccccc@{}}
\text{Pattern }n&1&2&3&4&5&10&n\\ \hline
\text{Dots}&5&9&13&17&21&41&4n+1
\end{array}
\]

\item An L-shape has two arms of \(n\) dots, but the corner dot is counted twice. Hence
\[
n+n-1=2n-1.
\]
\[
\begin{array}{@{}lccccccc@{}}
\text{Pattern }n&1&2&3&4&5&10&n\\ \hline
\text{Dots}&1&3&5&7&9&19&2n-1
\end{array}
\]
The \(-1\) comes from subtracting the corner dot, which is shared by the two arms.
\end{enumerate}

\paragraph{Q4. Using the \(n\)th terms.}
\begin{enumerate}[label=(\alph*),leftmargin=*,itemsep=4pt]
\item For the row of squares,
\[
3n+1=61,\qquad 3n=60,\qquad n=20.
\]
So pattern 20 uses 61 matchsticks.

\item For the tables,
\[
2n+2=30,\qquad 2n=28,\qquad n=14.
\]
So 14 tables are needed.

\item For the houses,
\[
5n+1=100,\qquad 5n=99,\qquad n=19.8.
\]
This is not a whole number, so a row of houses cannot use exactly 100 matchsticks.
\end{enumerate}

\paragraph{Q5. Drawing sequences.}
Any drawing with the correct number of objects is acceptable.
\begin{enumerate}[label=(\alph*),leftmargin=*,itemsep=4pt]
\item \(2n+3\): patterns 1, 2 and 3 have \(5,7,9\) objects. Start with 3 and add 2 each time.
\item \(4n-1\): patterns 1, 2 and 3 have \(3,7,11\) objects. Start with 3 and add 4 each time.
\item \(5n+2\): patterns 1, 2 and 3 have \(7,12,17\) objects. Start with 7 and add 5 each time.
\end{enumerate}

\section*{Part 2: Square numbers}

\paragraph{Q6. Squares of dots.}
\begin{enumerate}[label=(\alph*),leftmargin=*,itemsep=4pt]
\item A square of \(n\) by \(n\) dots has \(n^2\) dots.
\[
\begin{array}{@{}lccccccc@{}}
\text{Pattern }n&1&2&3&4&5&10&n\\ \hline
\text{Dots}&1&4&9&16&25&100&n^2
\end{array}
\]

\item Pattern 5 splits as
\[
25=1+3+5+7+9.
\]
The L-shape added to make pattern \(n\) has \(2n-1\) dots, because it has \(n\) dots on each arm and the corner dot is shared.
\end{enumerate}

\paragraph{Q7. Squares with dots added or taken away.}
\begin{enumerate}[label=(\alph*),leftmargin=*,itemsep=4pt]
\item A square of \(n^2\) dots plus 2 more gives \(n^2+2\).
\[
\begin{array}{@{}lccccccc@{}}
\text{Pattern }n&1&2&3&4&5&10&n\\ \hline
\text{Dots}&3&6&11&18&27&102&n^2+2\\
n^2&1&4&9&16&25&100&n^2
\end{array}
\]

\item Pattern \(n\) is a rectangle with \(n\) columns and \(n+1\) rows, so it has
\[
n(n+1)=n^2+n
\]
dots.
\[
\begin{array}{@{}lccccccc@{}}
\text{Pattern }n&1&2&3&4&5&10&n\\ \hline
\text{Dots}&2&6&12&20&30&110&n^2+n\\
n^2&1&4&9&16&25&100&n^2
\end{array}
\]

\item Two squares of \(n\) by \(n\) dots give
\[
2n^2.
\]
\[
\begin{array}{@{}lccccccc@{}}
\text{Pattern }n&1&2&3&4&5&10&n\\ \hline
\text{Dots}&2&8&18&32&50&200&2n^2
\end{array}
\]

\item A square of \(n^2\) dots with one dot removed gives
\[
n^2-1.
\]
\[
\begin{array}{@{}lccccccc@{}}
\text{Pattern }n&1&2&3&4&5&10&n\\ \hline
\text{Dots}&0&3&8&15&24&99&n^2-1
\end{array}
\]
\end{enumerate}

\paragraph{Q8. Staircases.}
\begin{enumerate}[label=(\alph*),leftmargin=*,itemsep=4pt]
\item The staircase counts are triangular numbers:
\[
1,\ 3,\ 6,\ 10,\ 15,\ 55.
\]

\item Two copies of pattern \(n\) fit together to make an \(n\) by \(n+1\) rectangle, so the rectangle has
\[
n(n+1)
\]
dots.

\item Therefore one staircase has half of this:
\[
\frac{n(n+1)}{2}.
\]
\[
\begin{array}{@{}lccccccc@{}}
\text{Pattern }n&1&2&3&4&5&10&n\\ \hline
\text{Dots}&1&3&6&10&15&55&\dfrac{n(n+1)}{2}
\end{array}
\]
\end{enumerate}

\paragraph{Q9. Drawing sequences.}
Any drawing with the correct number of objects is acceptable.
\begin{enumerate}[label=(\alph*),leftmargin=*,itemsep=4pt]
\item \(n^2+3\): patterns 1, 2 and 3 have \(4,7,12\) objects. Draw an \(n\) by \(n\) square and add 3 dots.
\item \(3n^2\): patterns 1, 2 and 3 have \(3,12,27\) objects. Draw three \(n\) by \(n\) squares.
\item \(n^2+2n=n(n+2)\): patterns 1, 2 and 3 have \(3,8,15\) objects. Draw an \(n\) by \(n+2\) rectangle of dots.
\end{enumerate}

\section*{Part 3: Pixel art sequences}

\paragraph{Q10. Pixel rocket.}
The rocket has 1 nose pixel, \(3n\) body pixels, a row of 5 pixels, and 2 extra side pixels. Hence
\[
1+3n+5+2=3n+8.
\]
\[
\begin{array}{@{}lccccccc@{}}
\text{Pattern }n&1&2&3&4&5&10&n\\ \hline
\text{Pixels}&11&14&17&20&23&38&3n+8
\end{array}
\]
\begin{enumerate}[label=(\alph*),leftmargin=*,itemsep=4pt,start=2]
\item For 50 pixels,
\[
3n+8=50,\qquad 3n=42,\qquad n=14.
\]
So rocket pattern 14 has 50 pixels.

\item For 100 pixels,
\[
3n+8=100,\qquad 3n=92,
\]
which does not give a whole number. So a rocket cannot have exactly 100 pixels.
\end{enumerate}

\paragraph{Q11. Pixel heart.}
Each of the 16 heart pixels becomes an \(n\) by \(n\) block, so the number of pixels is
\[
16n^2.
\]
\[
\begin{array}{@{}lccccccc@{}}
\text{Pattern }n&1&2&3&4&5&10&n\\ \hline
\text{Pixels}&16&64&144&256&400&1600&16n^2
\end{array}
\]
Pattern 2 is twice as wide and twice as tall as pattern 1. Each original pixel becomes a \(2\) by \(2\) block, so the number of pixels is multiplied by
\[
2\times 2=4,
\]
not by 2.

\paragraph{Q12. Pixel pyramid.}
The pyramid has rows of \(1,3,5,\ldots,2n-1\) pixels. These are the first \(n\) odd numbers, and their sum is
\[
1+3+5+\cdots +(2n-1)=n^2.
\]
\[
\begin{array}{@{}lccccccc@{}}
\text{Pattern }n&1&2&3&4&5&10&n\\ \hline
\text{Pixels}&1&4&9&16&25&100&n^2
\end{array}
\]
For pattern 4, moving the shaded pixels from the left side into the gaps on the right turns the pyramid into a \(4\) by \(4\) square. More generally, the pyramid rearranges into an \(n\) by \(n\) square, which has \(n^2\) pixels. Hence the pyramid has \(n^2\) pixels.

\paragraph{Q13. Frame around a square picture.}
The picture has \(n^2\) pixels. The frame is
\[
(n+2)^2-n^2=4n+4
\]
pixels. Hence the total is
\[
(n+2)^2=n^2+4n+4.
\]
\[
\begin{array}{@{}lccccccc@{}}
\text{Pattern }n&1&2&3&4&5&10&n\\ \hline
\text{Picture}&1&4&9&16&25&100&n^2\\
\text{Frame}&8&12&16&20&24&44&4n+4\\
\text{Total}&9&16&25&36&49&144&(n+2)^2
\end{array}
\]
The frame can be counted as a top row of \(n+2\) pixels, a bottom row of \(n+2\) pixels, and then \(n\) pixels on each of the left and right sides. Thus
\[
2(n+2)+2n=4n+4.
\]

\paragraph{Q14. Design your own sequence.}
One possible design is a robot with an \(n\) by \(n\) square body and one extra pixel at each corner. Pattern \(n\) then has
\[
n^2+4
\]
pixels.

\[
\begin{array}{@{}lccccccc@{}}
\text{Pattern }n&1&2&3&4&5&10&n\\ \hline
\text{Pixels}&5&8&13&20&29&104&n^2+4
\end{array}
\]
Many other designs are possible.

\section*{Part 4: Always, sometimes or never?}

\paragraph{Q15.} \textbf{Always.} \(3n+1\) is one more than \(3n\), and \(3n\) is a multiple of 3.

\paragraph{Q16.} \textbf{Never.} \(4n+2\) leaves remainder 2 when divided by 4, so it cannot be a multiple of 4. For example, when \(n=1\), \(4n+2=6\).

\paragraph{Q17.} \textbf{Sometimes.} If \(n=5\), then \(2n+5=15\), which is a multiple of 5. If \(n=1\), then \(2n+5=7\), which is not a multiple of 5.

\paragraph{Q18.} \textbf{Sometimes.} If the coefficient of \(n\) is even, such as \(2n+1\), every term is odd: \(3,5,7,\ldots\). If the coefficient is odd, such as \(3n+1\), the parity alternates: \(4,7,10,\ldots\), and \(4\) is even.

\paragraph{Q19.} \textbf{Sometimes.} For \(n=3\), \(n^2=9>6=2n\). But for \(n=1\), \(n^2=1\) is not greater than \(2\). Also \(n=2\) gives equality: \(4=4\).

\paragraph{Q20.} \textbf{Always.}
\[
n^2+n=n(n+1).
\]
These are consecutive integers, so one of them is even. Hence the product is even.

\paragraph{Q21.} \textbf{Never.} A square number ends in one of the digits \(0,1,4,5,6,9\). It cannot end in 7.

\paragraph{Q22.} \textbf{Never.} For \(n\ge 1\),
\[
n^2<n^2+1<(n+1)^2,
\]
so \(n^2+1\) lies between two consecutive square numbers and cannot itself be a square number.

\paragraph{Q23.} \textbf{Always.} The difference between consecutive square numbers is
\[
(n+1)^2-n^2=2n+1,
\]
which is odd.

\paragraph{Q24.} \textbf{Never.} The heart in pattern \(n\) has \(16n^2\) pixels. Doubling the pattern number gives
\[
16(2n)^2=64n^2=4(16n^2),
\]
which is four times as many pixels, not twice as many.

\paragraph{Q25.} \textbf{Always.} The frame has
\[
4n+4=4(n+1)
\]
pixels, which is a multiple of 4.

\paragraph{Q26.} \textbf{Sometimes.} A staircase in pattern \(n\) has
\[
\frac{n(n+1)}{2}
\]
dots. For \(n=1\), this is 1, which is odd. For \(n=3\), this is 6, which is even.

\paragraph{Q27. Own statements.}
One possible set of answers is:
\begin{itemize}
\item \textbf{Always:} A term of \(6n+2\) is even, because \(6n+2=2(3n+1)\).
\item \textbf{Sometimes:} A term of \(3n+2\) is a multiple of 5. This is true for \(n=1\), giving 5, but false for \(n=2\), giving 8.
\item \textbf{Never:} A term of \(5n+2\) is a multiple of 5, because it always leaves remainder 2 when divided by 5.
\end{itemize}

\end{document}